\begin{claim}{D.2.17}{Proved}{An actual realization on the six-skeleton}
Put $B=K(\FF_2,2)$, $X=\Sigma B$, $A=QX$, $Y=\Sigma^2B$, and
$T=\Omega^2Y$. We use the labels $a,b,c,d,e,f,x$ of
$\widetilde H_*(B;\FF_2)$ dual, respectively, to
\[
 u,\quad v=\operatorname{Sq}^1u,\quad u^2,\quad
 w=\operatorname{Sq}^2v,\quad uv,\quad u^3,\quad v^2.
 \tag{D.2.17a}
The same letters denote their images in $H_*T$. Write
$\xi=\xi_1$ for the top double-loop operation, and set
\[
 k=\{a,b\},\qquad p_3=f+ac+a^3,\qquad
 r=b^2+k+p_3\in H_6(T;\FF_2).
 \tag{D.2.17b}
There is an actual pointed map $\rho:S^6\to T$ with Hurewicz
image $r$.

There is a CW model of $A$ with finite six-skeleton $K$, and an
ordinary pointed map $G:K\to T$, with the following values on
the indicated lifts of the unit classes and the bottom product:
\[
\begin{array}{c|rrrrrr}
 z&A_0=\sigma a&B_0=\sigma b&C=\sigma c&D=\sigma d&E=\sigma e&A_0^2\\
 \hline
 G_*z&0&a^2&0&b^2+k&r&0.
\end{array}
 \tag{D.2.17c}
\]
Thus the remaining low-dimensional pattern with $s=1$, $\gamma=0$,
and $\epsilon=1$ has an actual six-dimensional realization.
This assertion does not give an extension over the seven-cells,
over $C_4X$, or over $A$.
\end{claim}
\begin{proof}
All unlabeled homology groups in the calculation have coefficients
$\FF_2$. We write $\beta=\operatorname{Sq}_*^1$; the notation
$\beta_2$ always means the differential on the second page of the
integral Bockstein spectral sequence.

\medskip\noindent
\textbf{The degree-six primitive and its second Bockstein.}
The free double-loop homology basis and the coproduct give
\[
 \operatorname{Prim}H_6T=\langle b^2,k,p_3\rangle,
 \qquad
 \operatorname{Sq}_*^2(b^2)=\operatorname{Sq}_*^2(p_3)=a^2,
 \qquad \operatorname{Sq}_*^2(k)=0.
 \tag{D.2.17d}
\]
For completeness, the entire degree-six basis is
$b^2,k,f,x,ac,a^3$. The class $x$ has the unique $b\otimes b$
term in its reduced coproduct. The remaining nonprimitive terms are
\[
 \overline\Delta f=a\otimes c+c\otimes a,
 \qquad
 \overline\Delta(ac)=\overline\Delta f+a^2\otimes a+a\otimes a^2,
 \qquad
 \overline\Delta(a^3)=a^2\otimes a+a\otimes a^2.
\]
This proves the asserted primitive basis. Cartan and the label
operations give the square calculation. These are the conventions
of \cite[Chapter~III, Theorems~1.2--1.4 and 3.1--3.2]{I_CohenLadaMay1976}.

We need the following integral refinement:
\[
 [b^2]=[k],\qquad
 \beta_2[b^2]=\beta_2[k]=[\xi a]\ne0,
 \qquad [p_3]=0
 \quad\hbox{on the second Bockstein page.}
 \tag{D.2.17e}
\]
Here is a chain calculation fixing its coefficient. Use the bottom
Moore map $M(\ZZ/2,2)\to B$, with integral reduced chains $a_2,b_3$
and $db=2a$. In the arity-two configuration circle take free
$\Sigma_2$-cells $e_0,e_1$ with $de_1=(\tau-1)e_0$.
The action on tensors includes the Koszul sign:
$\tau(a\otimes b)=b\otimes a$ and
$\tau(b\otimes b)=-b\otimes b$. In the coinvariant tensor complex put
\[
 z=e_1(a\otimes a),\qquad
 q=e_0(b\otimes b),\qquad h=e_1(a\otimes b).
\]
With the total-complex differential,
\[
\begin{split}
 dq&=2\bigl(e_0(a\otimes b)-e_0(b\otimes a)\bigr),\\
 dh&=e_0(b\otimes a)-e_0(a\otimes b)-2z,\qquad dz=0,\\
 d(q+2h)&=-4z.
\end{split}
 \tag{D.2.17f}
\]
The reductions of $q$ and $z$ represent $b^2$ and $\xi a$.
Thus the corrected integral representative $q+2h$ proves
$\beta_2[b^2]=[\xi a]$. In particular, we are not asserting that
$dq$ is divisible by four. The same model gives
$\beta(\xi b)=b^2+k$. The configuration-chain conventions are those
in the proof of \cite[Chapter~III, Lemma~3.10, p.~232]{I_CohenLadaMay1976};
naturality along the bottom Moore map gives these identities in $T$.
There is no degree-six first-Bockstein preimage of $\xi a$:
on the complete degree-six basis its only nonzero value is
$\beta x=d+e$. Finally, with $g$ the label dual to $u^2v$,
\[
 \beta(g+ae+bc+a^2b)=p_3.
\]
This proves all parts of \textup{(D.2.17e)}.

It follows that a spherical degree-six class with nonzero double
homology suspension must be $r$. Indeed a spherical class is primitive
and is killed by $\operatorname{Sq}_*^2$ and $\beta_2$. The square
condition in \textup{(D.2.17d)} makes its coefficients on $b^2$ and
$p_3$ equal, and the Bockstein condition makes its coefficients on
$b^2$ and $k$ equal. Double homology suspension kills $b^2$ and $k$
and sends $p_3$ to $\sigma^2f$.

\medskip\noindent
\textbf{The required class is actually spherical.}
Let $M=\mathbb RP^3\times\mathbb RP^3$.
The proof of \cite[Theorem~1.1, p.~9]{EH5AkhmetievRepovsSkopenkov2001}
constructs a smooth embedding
\[
 M\times I^2\hookrightarrow\mathbb R^8.
 \tag{D.2.17g}
\]
Restricting to the middle slice gives an embedding of $M$ with a
trivialized rank-two normal bundle: the two interval directions
project to a frame of its normal quotient. Write
$t_1,t_2\in H^1(M;\FF_2)$ for the two projective-space generators,
and choose $j:M\to B$ classifying $t_1t_2$. The framed
Pontryagin--Thom collapse gives an actual map
\[
 \theta:S^8\longrightarrow
 \operatorname{Th}(\varepsilon_M^2)
 \cong S^2\wedge M_+
 \xrightarrow{S^2\wedge j_+} S^2\wedge B=Y,
 \tag{D.2.17h}
\]
where $j_+$ sends the disjoint basepoint to the basepoint of $B$.
Its Hurewicz image is $\sigma^2j_*[M]$. Serre's calculation gives
the degree-six cohomology basis $u^3,v^2$
\cite[\S2, nos.~8--9, pp.~203--204]{I_Serre1953}, and
\[
 \langle j^*(u^3),[M]\rangle
   =\langle t_1^3t_2^3,[M]\rangle=1,
 \qquad
 j^*(v^2)=(t_1^2t_2+t_1t_2^2)^2
          =t_1^4t_2^2+t_1^2t_2^4=0.
 \tag{D.2.17i}
\]
Hence $j_*[M]=f$. The double adjunct $\rho:S^6\to T$ has double
homology suspension $\sigma^2f$, so the preceding uniqueness
calculation gives $\rho_*[S^6]=r$. This uses the actual rank-two
normal framing in \textup{(D.2.17g)}, rather than a framing in a
larger Euclidean space.

\medskip\noindent
\textbf{The canonical map before correction.}
Let
\[
 h_0=\eta_2\wedge\id_B:\Sigma^3B\to\Sigma^2B,
 \qquad
 g_2=\Omega^2h_0:C_2X\to T.
 \tag{D.2.17j}
\]
On the unit $X\to C_2X$ its restriction $g_0$ has values
\[
 g_0(A_0)=0,\quad g_0(B_0)=a^2,\quad g_0(C)=\xi a,
 \quad g_0(D)=b^2,\quad g_0(E)=k.
 \tag{D.2.17k}
\]
We include the check of the entry on $B_0$, which is not a formal
consequence of the entry on $C$. Write $X_7=\sigma x$.
The cofiber of $h_0$ is $\mathbb CP^2\wedge B$. Its reduced
diagonal gives the quadratic once-adjoint image of $X_7$ as the
square of the loop label $\sigma b$. The linear part is zero,
since $h_0$ is zero on homology; tensor length three begins in
degree nine, above this degree-eight calculation. This is the
relative Adams--Hilton computation used in D.2.13, with the
quadratic differential and cell-attachment conventions of
\cite[\S3.1, p.~300; proof of Proposition~7.4, p.~309]{FGHHalperinLemaire1995}.

The primitive degree-seven basis of $H_*T$ is
\[
 \xi b,\quad \{a,c\},\quad
 P=g+ae+bc+a^2b,\quad Q=h+ad,
 \tag{D.2.17l}
\]
where $g,h$ are the labels dual to $u^2v,uw$.
Their homology suspensions are, respectively, the square of
$\sigma b$, the commutator of $\sigma a$ and $\sigma c$,
and the two independent linear labels $\sigma g,\sigma h$.
Thus homology suspension is injective on this primitive subspace,
and $g_0(X_7)=\xi b$. The top operation formula gives
\[
 \operatorname{Sq}_*^3(\xi b)=(\beta b)^2=a^2,
 \qquad \operatorname{Sq}_*^3X_7=B_0.
\]
Naturality proves $g_0(B_0)=a^2$. The value on $C$ was proved in
D.2.13. The degree-three primitive basis is $\langle b\rangle$,
and $\beta b=a\ne0$, whereas $\beta A_0=0$; hence $g_0(A_0)=0$.
The same quadratic diagonal calculation has zero value on $D$
and the commutator value on $E$. By \textup{(D.2.17d)}, the
kernel of homology suspension on the degree-six primitives is
$\langle b^2\rangle$. Finally
$\operatorname{Sq}_*^2D=B_0$ and
$\operatorname{Sq}_*^2E=0$ determine the possible $b^2$
coefficients, giving the remaining two entries of
\textup{(D.2.17k)}. Since $g_2$ is a double-loop map, it also sends
the product $A_0^2$ to zero.

\medskip\noindent
\textbf{The quotient of the six-skeleton.}
The map $C_2X\to QX$ is seven-connected. Indeed $\Sigma^2X$ is
four-connected, its Freudenthal unit is nine-connected, and two
loops leave connectivity seven; successive stabilization maps
have increasing connectivity. A relative CW approximation thus
adds no cells below dimension eight. We choose finite skeleta of
$C_2X$ and the relative approximation, and take $K=A^{(6)}$ from
this model. Its four-skeleton may be taken to be
$K^{(4)}=M(\ZZ/2,3)$. The unit classes have cellular representatives
in $K$, and the bottom product is represented in the second
configuration stage.

After collapsing $K^{(4)}$, the quotient is four-connected and
has cells only in dimensions five and six. Moreover
\[
 H_5(K/K^{(4)};\ZZ)=H_5(A;\ZZ)=H_4(B;\ZZ)=\ZZ/4.
 \tag{D.2.17m}
\]
Here $H_4(K^{(4)};\ZZ)=0$, and cells of dimension at least seven
do not affect $H_5$. The stable splitting
\cite[Theorem~1.1 and \S2]{Snaith1974} identifies $H_5A$ with
the weight-one group: all weights at least two begin in dimension
six for $X=\Sigma B$. The last equality follows from the
order-four calculation
\cite[proof in \S6, p.~1340]{I_Mahowald1964}, which identifies
$H^5(B;\ZZ)=\ZZ/4$, and the integral universal coefficient theorem;
the relevant homology is finite two-primary torsion.

The six-cell attachments to a wedge of five-spheres are entirely
determined by their degree matrix. Smith normal form and cell
slides therefore give
\[
 K/K^{(4)}\simeq M(\ZZ/4,5)\vee\bigvee_j S^6,
 \tag{D.2.17n}
\]
after removal of contractible cell pairs. We retain the free
six-spheres. Their role is essential in the following choice.
The order-four Bockstein has
$\beta_2[D]=\beta_2[E]=[C]$. Equivalently, the integral
Pontryagin-square generator has reduction $w+uv$
\cite[\S2.3, p.~1336]{I_Mahowald1964}.
Thus, in an integral cellular choice lifting these mod-two
classes, the boundary degrees of the cells $D,E$ on the
five-cell are $4,-4$, respectively. Their sum is an integral
free six-cycle. The bottom product $A_0^2$ is an independent
free six-cycle, in the configuration weight-two summand.

We can consequently choose pointed maps
\[
 q:K/K^{(4)}\to M(\ZZ/4,5),\qquad
 \ell:K/K^{(4)}\to S^6
 \tag{D.2.17o}
\]
with the following values. The map $q$ has degree one on the
five-cell, sends $D,E$ to the mod-two top Moore class, and
kills $A_0^2$. For $\ell$, collapse the five-skeleton and project
to the $E$ six-cell. Then
\[
 \ell_*D=0,\qquad \ell_*E=[S^6],\qquad \ell_*(A_0^2)=0.
 \tag{D.2.17p}
\]
Any additional free six-spheres are sent to the basepoint.
These maps are choices on the actual quotient cell complex.
No equation involving a seven-dimensional source class is being
imposed on an arbitrary map from $K$.

\medskip\noindent
\textbf{The Moore correction and its top-cell choice.}
Let $u_4:S^4\to Y$ be the bottom map, of order two by Hurewicz,
and let $\nu_4:S^7\to S^4$ be the quaternionic Hopf map.
Set $\alpha=u_4\nu_4$. The Hopf--Hilton formula gives
\[
 0=(2u_4)\nu_4=2\alpha+[u_4,u_4],
 \qquad 2[u_4,u_4]=[2u_4,u_4]=0,
 \tag{D.2.17q}
\]
so $4\alpha=0$. The formula is obtained by pinching the target
of $\nu_4$ into two four-spheres: the only additional component
in this dimension is the Whitehead product, with coefficient
the Hopf invariant one. Its quadratic form and composition rule
are recorded in \cite[(8.2)(4), (8.4), and Lemma~8.5,
pp.~4756--4757]{I_Baues1996}.

The double adjunct $\overline\alpha:S^5\to T$ has Hurewicz
image $\xi a$. To verify the coefficient, first work with
$\nu_4$: its cofiber is $\mathbb HP^2$, whose bottom class has
nonzero square. Its once-adjoint Hurewicz image is therefore
the square of the degree-three loop generator. Homology
suspension in $\Omega^2S^4$ sends $\xi a$ to that square and
has zero primitive kernel in degree five. Naturality under
$u_4$ proves the assertion in $T$.

Since $4\overline\alpha=0$, it extends to a map
$\mu:M(\ZZ/4,5)\to T$. Its top mod-two class is primitive,
is killed by $\operatorname{Sq}_*^2$, and has second
Bockstein $\xi a$. Equations \textup{(D.2.17d)--(D.2.17e)}
therefore force its image to have the form
\[
 k+t r,\qquad t\in\FF_2.
 \tag{D.2.17r}
\]
The top-cell coaction
$M(\ZZ/4,5)\to M(\ZZ/4,5)\vee S^6$, followed by
$\mu\vee\rho$, changes $t$ by one while retaining the given
bottom map. We choose the extension with top image exactly $k$.

Let $p:K\to K/K^{(4)}$ be the quotient map and form the
pointwise sum in the H-group $T$:
\[
 G=g_2|_K+\mu q p+\rho\ell p:K\longrightarrow T.
 \tag{D.2.17s}
\]
Both corrections vanish on $K^{(4)}$. On $C$ the first
correction adds $\xi a$ and the second adds zero. On $D,E$,
the first adds $k$ to each, while the second adds zero to $D$
and $r$ to $E$. All these source classes are primitive, so
the homology of the pointwise sum has no cross terms. The
product class $A_0^2$ is also primitive in characteristic two;
$g_2$ kills it by multiplicativity and both corrections kill
its chosen free-sphere representative. Combining this with
\textup{(D.2.17k)} gives exactly \textup{(D.2.17c)}.

The separate free-sphere choice in \textup{(D.2.17p)} is what
sets the value on $E$ correctly. The construction does not
verify the seven-cell attaching relations or the required
integral orders of their images. Those are further extension
conditions, and neither a map on $A$ nor an ordinary
nullification calculation is inferred from the finite map.
\end{proof}
